\chapter{Models}\label{chap:models}

\section{Baseline Models}\label{sec:models-baselines}

Baseline models provide reference performance for evaluating proposed
architectures. Candidate baselines include classical machine learning models,
CNNProm-style convolutional networks, GraphPro-inspired approaches, and
transformer-based promoter models where feasible. Neural-network notation and
model descriptions follow the distinction between data representations, model
functions, trainable parameters, and objective functions commonly used in
deep-learning methodology~\citep{goodfellow2016deep}. In this thesis, DNA-BERT has
two planned roles: it may provide sequence embeddings for other classifiers, and
it may also be fine-tuned as an end-to-end promoter classifier with a binary
classification head. The latter role belongs to the model-development stage and
is intentionally kept separate from the feature-engineering encoder.

\section{CapsProm}\label{sec:models-capsprom}

CapsProm is the capsule-network approach developed in the earlier stage of this
research. In this thesis, CapsProm should be reproduced under the unified
experimental protocol before architectural extensions are evaluated.

\section{SparseCapsProm}\label{sec:models-sparsecapsprom}

SparseCapsProm is reserved for capsule-network variants designed to reduce
model complexity or encourage sparse feature interactions. \pendingresults

\section{Attention-Enhanced CapsProm}\label{sec:models-acapsprom}

The planned A-CapsProm architecture will combine capsule-based modeling with
attention mechanisms. This extension should be evaluated only after baseline and
CapsProm reproduction experiments are stable.
